% ==================================================================
% §15.5.6 REBUILT — Phase 1: Interpretation and outlook
%
% Part 6 (FINAL) of the §15.5 rebuild. Insert IMMEDIATELY AFTER
% §15.5.5 block (label: subsec:eps_comparison_rebuilt) and BEFORE
% the existing old §15.5.
%
% GPT round 13 (6 corrections) applied:
%   (1) Removed "without fine-tuning" overclaim; added "not a full
%       global likelihood result" qualifier; "qualitatively distinct".
%   (2) Sharpened non-independence: "not a measurement of a first-
%       principles constant and not a statistically rigorous global
%       combination of independent likelihoods".
%   (3) OP-Hubble-derive paragraph softened: "For reference within
%       the present rebuilt §15.5, we denote this task by
%       OP-Hubble-derive" (not yet in central registry).
%   (4) "Parsimony and correlations" softened: closure-level
%       extension framed as structural expectation only; quantitative
%       realisation belongs to OP-Hubble-derive.
%   (5) "Three layers" sharpened with closing: "No statement in the
%       rebuilt §15.5 should be read as belonging to the first-
%       principles closure level unless explicitly said otherwise".
%   (6) Final meta-bridge "This closes the rebuilt §15.5." removed.
%
% Absorbs inventory items: #6, #7, #8 qualitative, #12, #14a.
%
% NO legacy numbers. NO new table, no new figure.
% ==================================================================

\subsubsection*{\normalfont\textbf{15.5.6 \quad Interpretation and outlook}}
\label{subsec:eps_outlook_rebuilt}

With the retained five-probe effective band of
\S\ref{subsec:eps_joint_band_rebuilt} and its methodological context of
\S\S\ref{subsec:eps_retained_rebuilt}--\ref{subsec:eps_comparison_rebuilt}
in place, we close this subsection by stating what the analysis
establishes, what it does not, and where the ECT $\varepsilon$-sector
programme stands.

\paragraph{What the analysis establishes.}
Under the stated assumptions (the uniform-$\varepsilon$ ansatz of
\S\ref{subsec:eps_def_rebuilt}, the $\Lambda$CDM-background proxy, the
physical prior $\varepsilon \geq 0$, and the retention criteria of
\S\ref{subsec:eps_retained_rebuilt}), the five retained probes admit a
narrow joint effective allowed band. ECT's three-stage
condensate-closure picture is compatible with this result at the
diagnostic-layer level. The compatibility of a single effective
deformation parameter with five channels of qualitatively distinct
observational systematics constitutes an indirect consistency check
of the framework, not a derivation from first principles and not a
full global likelihood result.

\paragraph{What the analysis does NOT establish.}
It does not establish $\varepsilon$ as a fundamental parameter of
ECT. The retained band is empirical under the uniform-$\varepsilon$
diagnostic layer, not a measurement of a first-principles constant
and not a statistically rigorous global combination of independent
likelihoods. It does not provide a completed resolution of the
Hubble or JWST tensions in the sense of a full CMB-to-local Boltzmann
pipeline; the present treatment implements the late-time H1
closure-level route, not the full H2 CMB-inferred inference pipeline.
It does not definitively exclude the inadequacy of a single effective
uniform-$\varepsilon$; broader datasets or methodological improvements
on currently excluded channels might reveal tensions requiring an
epoch-dependent $\varepsilon(z)$.

\paragraph{Diagnostic layer, not final closure.}
The uniform-$\varepsilon$ treatment is used here as a diagnostic
effective layer, not as the final closure-level description of the
cosmological drift sector. A full ECT-native treatment of the
cosmological drift --- including the functional form of
$\varepsilon(z)$ implied by three-stage condensate closure ---
remains open.

\paragraph{Open problem: $\varepsilon(z)$ from closure.}
\label{op:hubble_derive}
A first-principles ECT derivation of the cosmological drift
parameter --- specifically, the closure-level function
$\varepsilon(z)$ implied by the three-stage condensate closure,
together with its mapping to the effective uniform-$\varepsilon$
analysis layer of \S\ref{subsec:eps_def_rebuilt} --- remains open.
For reference within the present rebuilt \S15.5, we denote this task
by \textbf{OP-Hubble-derive}.

\paragraph{Methodological requirements.}
Three methodological improvements are noted here without being logged
as formal open problems, since each is a question of extending an
analysis pipeline rather than a structural theoretical gap. The BAO
channel of \S\ref{subsec:eps_excluded_rebuilt} requires a full DESI
likelihood with a coupled $(\Omega_m, r_d, \varepsilon)$ fit before it
can be considered for the joint band at headline level. The
$A_{\rm lens}$ channel of \S\ref{subsec:eps_excluded_rebuilt} requires
a full Boltzmann-level implementation before it can be similarly
considered. A genuine SN~Ia $\varepsilon$-channel requires a real
observational supernova pipeline with full SH0ES and host-galaxy
systematics.

\paragraph{Parsimony and correlations.}
In a closure-level extension, the same underlying drift mechanism
would be expected to correlate the late-time expansion history, the
effective dark-energy phenomenology, and the time budget available
for early structure formation. This qualitative correlation structure
is a distinctive feature of the ECT use of the $\varepsilon$-sector
rather than a post-hoc collection of separate fitting knobs. At the
present stage, however, this remains only a structural expectation:
its quantitative realisation belongs to OP-Hubble-derive.

\paragraph{Three layers of the analysis.}
Three distinct layers coexist in the present treatment and should be
kept separate: the ECT structural motivation for a cosmological
deformation of the gravitational response
(\S\ref{subsec:eps_def_rebuilt}); the uniform-$\varepsilon$ diagnostic
layer adopted in \S15.5 and constrained empirically by the retained
five-probe band
(\S\S\ref{subsec:eps_retained_rebuilt}--\ref{subsec:eps_joint_band_rebuilt});
and the closure-level $\varepsilon(z)$ derivation that remains open
(OP-Hubble-derive). No statement in the rebuilt \S15.5 should be read
as belonging to the first-principles closure level unless explicitly
said otherwise; all quantitative conclusions of \S15.5 belong to the
middle diagnostic layer only.

\paragraph{Direction for future work.}
Three natural directions follow. First, methodological upgrades to
currently excluded channels (BAO, $A_{\rm lens}$, SN~Ia) would test
whether the retained band remains stable as the analysis base
broadens. Second, a proper $\varepsilon(z)$ parameterisation, derived
from closure theory when it matures, would test whether the
uniform-$\varepsilon$ assumption is the adequate shape of the
deformation in this epoch window. Third, a systematic study of the
dependence of the retained band on the $\Lambda$CDM-background proxy
would quantify the residual proxy-dependence of the present
methodology.

% ==================================================================
% END §15.5.6 (rebuilt). All six rebuilt subsections §15.5.1--§15.5.6
% are now in place. DO NOT REMOVE existing §15.5 below until Phase 3
% switchover per PROJECT_STRUCTURE §3.6.
% ==================================================================
